\subsection{COXETER MATRIX AND COXETER GRAPH OF W}

Let C be a chamber, \(S = S(C)\) the set of orthogonal reflections with respect to the walls of C and \(M = (m(s, s'))\) the Coxeter matrix of the Coxeter system (W, S) (Chap. IV, §1, no. 9): recall that \(m(s, s')\) is the order (finite or infinite) of the element \(ss'\) of W (for \(s, s' \in S\) ). If \(C'\) is another chamber, the unique element \(w \in W\) such that \(w(C) = C'\) defines a bijection

\[
s \mapsto f (s) = w s w ^ {- 1}
\]

from S to \(S' = S(C')\) , and we have \(m(f(s), f(s')) = m(s, s')\) . It follows that, if W acts on the set X of pairs \((C, s)\) , where C is a chamber and \(s \in S(C)\) , by \(w.(C, s) = (w(C), wsw^{-1})\) , each orbit i of W in X meets each of the sets \(\{C\} \times S(C)\) in exactly one point, which we denote by \((C, s_i(C))\) . Thus, if I is the set of orbits and \(i, j \in I\) , the number \(m_{ij} = m(s_i(C), s_j(C))\) is independent of the choice of chamber C. The matrix \(M(W) = (m_{ij})_{i,j \in I}\) is a Coxeter matrix called the Coxeter matrix of W. The Coxeter graph associated to M(W) (Chap. IV, §1, no. 9) is called the Coxeter graph of W.

Let C be a chamber. For any \(i \in I\) , denote by \(\mathrm{H}_{i}(\mathrm{C})\) the wall of C such that \(s_{i}(\mathrm{C})\) is the reflection with respect to \(\mathrm{H}_{i}(\mathrm{C})\) and by \(e_{i}(\mathrm{C})\) the unit vector orthogonal to \(\mathrm{H}_{i}(\mathrm{C})\) on the same side of \(\mathrm{H}_{i}(\mathrm{C})\) as C. The map \(i \mapsto \mathrm{H}_{i}(\mathrm{C})\) is called the canonically indexed family of walls of C.

PROPOSITION 3. Let \(\mathbf{C}\) be a chamber and let \(i, j \in \mathbf{I}\) with \(i \neq j\) . Put \(s_i = s_i(\mathbf{C})\) , \(\mathbf{H}_i = \mathbf{H}_i(\mathbf{C})\) , \(e_i = e_i(\mathbf{C})\) and define \(s_j\) , \(\mathbf{H}_j\) and \(e_j\) similarly.

\begin{enumerate}
\def\labelenumi{(\roman{enumi})}
\item
  If \(\mathrm{H}_i\) and \(\mathrm{H}_j\) are parallel, then \(m_{ij} = \infty\) and \(e_i = -e_j\) .
\item
  If \(\mathrm{H}_i\) and \(\mathrm{H}_j\) are not parallel, then \(m_{ij}\) is finite and
\end{enumerate}

\[
\left(e _ {i} \mid e _ {j}\right) = - \cos (\pi / m _ {i j}).\tag{4}
\]

\begin{enumerate}
\def\labelenumi{(\roman{enumi})}
\setcounter{enumi}{2}
\tightlist
\item
  We have \((e_i|e_j)\leqslant 0\)
\end{enumerate}

If \(\mathrm{H}_i\) and \(\mathrm{H}_j\) are parallel, \(s_i s_j\) is a translation (§2, no. 4, Prop. 5), so \(m_{ij} = \infty\) . Moreover, either \(e_i = e_j\) or \(e_i = -e_j\) . Now, there exists a point \(a\) (resp. \(a'\) ) in the closure of C that belongs to \(\mathrm{H}_i\) (resp. \(\mathrm{H}_j\) ) but not to \(\mathrm{H}_j\) \((\text{resp. } \mathrm{H}_i)\) . Then \((a' - a|e_i) > 0\) and \((a - a'|e_j) > 0\) , which excludes the case \(e_i = e_j\) and proves (i).

Assume now that \(H_{i}\) and \(H_{j}\) are not parallel. Choose an origin \(a \in H_{i} \cap H_{j}\) and identify T with E by the bijection \(t \mapsto a + t\) . Let V be the plane orthogonal to \(H_{i} \cap H_{j}\) and passing through a. Put \(\Gamma = V \cap D_{H_{i}}(C) \cap D_{H_{j}}(C)\) (where \(D_{H}(C)\) denotes the open half-space bounded by H and containing C (§1, no. 1)) and let D (resp. \(D'\) ) be the half-line in V contained in \(H_{i} \cap V\) (resp. \(H_{j} \cap V\) ) and in the closure of \(\Gamma\) . For a suitable orientation of V, the set \(\Gamma\) is the union of the open half-lines \(\Delta\) in V such that

\[
0 <   (\widehat {\mathrm{D} , \Delta}) <   (\widehat {\mathrm{D} , \mathrm{D} ^ {\prime}}).
\]

Let \(W'\) be the subgroup of \(W\) generated by \(s_i\) and \(s_j\) . For any \(w \in W'\) , the hyperplanes \(w(H_i)\) and \(w(H_j)\) belong to \(\mathfrak{H}\) , contain \(H_i \cap H_j\) and do not meet C. It follows that they do not meet \(\Gamma\) (§1, no. 5, Prop. 10). The Cor. of Prop. 7 of §2, no. 5 thus implies (ii).

Finally, assertion (iii) follows immediately from (i) and (ii), since \(m_{ij} \geqslant 2\) for \(i \neq j\) .

Remark. Formula (4) is actually valid for all \(i, j \in I\) : in fact, \(\pi / m_{ij} = 0\) if \(m_{ij} = \infty\) , and if \(i = j\) then \(m_{ij} = 1\) and \((e_i | e_j) = 1\) .

